\section{Limitations and conclusion}
\label{sec:conclusion}

\paragraph{Limitations.} The kernel analysis is exact at infinite width and
for the squared loss; the Bradley--Terry and feature-learning results are
empirical. Our networks are small MLPs on low-dimensional features, not reward
models on top of a language model, where the ``items'' are responses with rich
embeddings and each response appears in few comparisons --- a sparse regime in
which our ratio is smallest. The synthetic truth mixes one utility with one
skew interaction; other cyclic structures could be learned at different
speeds. The Arena analysis is at model level, so it says nothing about
prompt-level intransitivity, where genuine cycles are more plausible, and
held-out votes are split within pairs, so they test fitting of observed pairs
more than generalisation to unseen ones. The NTK stopping rule we proposed
failed; a rule that estimates the signal fraction in cycle space, for example
with the transitive null of Section~\ref{sec:arena}, is the natural
replacement. The full-scale configuration widens every sweep.

\paragraph{Conclusion.} Gradient descent learns preferences in a fixed order:
the ranking first, the cycles much later. For a scalar reward model the
cycles are not merely unrepresentable but invisible to the gradient. For an
antisymmetric network the edge NTK puts transitive flows 4 to over 100 times
higher in the spectrum than cyclic ones, more so on denser comparison graphs,
and training follows that ordering, closely in the lazy regime and robustly
beyond it. Local and global cycles, by contrast, are learned at the same
rate. The ordering turns training time into a knob on transitivity. It
explains why early stopping helps cyclic preference models most when votes are
noisy. It also cautions against reading learned cycles as signal: in Chatbot
Arena, the cycles a flexible model eventually learns are sampling noise,
fitted last.
